\section{Results}
\label{sec:results}

\begin{center}
\fbox{\parbox{0.88\linewidth}{%
\small
Sample note. Results below are reported on the \emph{25-country} sample.
The pre-specified benchmark covered 15 countries and was extended post hoc by ten
(six Global North references --- the United Kingdom, Canada, Australia, South
Korea, France, Italy --- and four native-pair countries --- Colombia, Chile,
Portugal, Angola). Fourteen LLMs are scored across five tasks, two personas, and
four languages (English plus a country-native language --- Spanish, Portuguese,
or Hindi --- for the nine applicable countries, giving the within-country
English/native pairs that test H2). The primary family recomputed on the
pre-specified 15 countries alone is tabulated in Supplementary Material, and no
pre-specified conclusion is revised on the extension alone; where the two samples
differ on a hypothesis the paper turns on, we say so in the text. Scoring is
adjudicated by code wherever the answer is a number checkable against an official
register, and by the mean of a three-vendor judge panel elsewhere
(Section~\ref{sec:results:reliability}); the T1 ground truth is anchored to the
national instrument in force for 24 countries and has no key for Egypt, which is
excluded from T1 (Section~\ref{sec:results:gt}).
}}
\end{center}

\subsection{Sample and scoring}
\label{sec:results:sample}

The analysed corpus comprises 8{,}244 scored cells on the $[0,1]$ composite,
of which 6{,}573 are English-prompt responses across the 25 countries; the
remainder are the within-country native-language pairs used for H2. Scoring uses
the most reliable instrument each task admits: code adjudicates against the
official register wherever the answer is a published number, resolving 99.8\% of
T1, 56.3\% of T2 and 71.8\% of T3 cells, and the residual, with all of T4, is
scored by the mean of a three-vendor panel. T5 retains the scores of the
collection judge. In total 6{,}434 cells carry a code or panel score.

The panel mean replaces the single-judge design of the original protocol, in which
the collection judge (GPT-5-mini for the pre-specified fifteen countries, Claude
Haiku~4.5 for the ten extension countries) was itself among the evaluated
models and, because the extension holds seven of the ten Global North countries,
was confounded with tier. Judges differ sharply
in severity (mean composite $0.617$ for Gemini, $0.451$ for Claude, $0.379$ for
DeepSeek), so a single judge fixes an arbitrary level, and single-judge reliability
is only ICC$(2,1)=0.56$ against ICC$(2,3)=0.79$ for the panel mean. One panel
member (DeepSeek-V3) is also an evaluated model, so we checked for self-favouring:
it scores its own outputs \emph{more} harshly relative to the other two judges
($-0.177$) than it scores any other model's ($-0.123$ to $-0.160$).

\subsection{Accuracy by model}
\label{sec:results:model}

Table~\ref{tab:conf-model} orders the 14 models by mean composite: frontier systems
lead and small open-weight models trail. On H3, the Brazilian-Portuguese regional
model Cabra-Mistral~7B is the \emph{weakest} of all ($0.351$ versus $0.557$ for the rest; Mann-Whitney $p\approx 0$, Cliff's
$\delta=-0.48$). Cliff's $\delta$ has a direct reading --- the probability that a
randomly drawn answer from one group beats one from the other, rescaled to
$[-1,+1]$ --- so $-0.48$ means the regional model loses roughly three such
comparisons for every one it wins. That comparison is against the whole roster,
which includes frontier systems; the pre-specified test was narrower, Cabra
against its scale-matched peer Llama~3.1~8B on Portuguese-language prompts, and
it points the same way with a smaller margin: on Brazil the regional model scores
$0.258$ against $0.313$ ($\delta=-0.26$, $n=20$ each), and across the three
Portuguese-speaking countries $0.304$ against $0.362$ ($\delta=-0.17$, $n=60$
each). Regional instruction tuning at the 7B level does not recover what scale
and frontier training give, even in the language and country it was tuned for,
though with 20 Brazilian cells per model the narrow test is descriptive.

\begin{table}[h]
\centering
\caption{Composite accuracy by model across 25 countries (code-adjudicated where
the answer is a register value, panel mean elsewhere; $[0,1]$). $N$ = scored
English-prompt responses.}
\label{tab:conf-model}
\footnotesize
\setlength{\tabcolsep}{8pt}
\renewcommand{\arraystretch}{1.15}
\begin{tabular}{lrr}
\toprule
Model & $N$ & Mean \\
\midrule
DeepSeek-V3      & 496 & 0.690 \\
GPT-5            & 496 & 0.676 \\
Gemini~2.5~Flash & 486 & 0.636 \\
GPT-5-mini       & 495 & 0.628 \\
Claude Haiku~4.5 & 486 & 0.589 \\
Qwen3 32B        & 494 & 0.560 \\
Command~R+       & 496 & 0.530 \\
Llama~3.3 70B    & 391 & 0.514 \\
Llama~4 Scout    & 496 & 0.500 \\
Qwen3 14B        & 496 & 0.485 \\
GPT-OSS 120B     & 257 & 0.479 \\
Phi-4 14B        & 496 & 0.475 \\
Llama~3.1 8B     & 496 & 0.437 \\
Cabra-Mistral~7B & 492 & 0.351 \\
\bottomrule
\end{tabular}
\end{table}

\subsection{Accuracy by task: the recall floor}
\label{sec:results:task}

Difficulty is sharply task-dependent (Table~\ref{tab:conf-task}). The hardest
task by a wide margin is T1 (technical-standard recall, 0.394) --- the
task that requires returning the national annual PM\textsubscript{2.5} standard
in force. The open recommendation task (T5) scores highest (0.769) because its
rubric rewards plausible, well-structured policy advice that does not hinge on a
single fact. Pooling the two recall tasks (T1, T2; mean $0.435$) against
synthesis and recommendation (T3--T5, $0.612$), the difference is large
(Mann-Whitney $p\approx 0$, Cliff's $\delta=-0.41$, roughly two and a half
losses for every win). LLMs are weakest exactly where applied environmental management
needs them most: recalling the precise
regulatory threshold in force.

The two corrections of the scoring are themselves informative here. T2 was
originally scored against a placeholder key and rises from $0.368$ to $0.473$
once each returned concentration is compared with the authorised range by code:
the correction moved the task whose key was defective. T1 always had an official
key but was scored by two collection judges of different severity, and moving it
to code raises it from $0.367$ to $0.394$ while keeping it the hardest task;
the code verdict also shows what the judges could not, that a third of the
wrong T1 answers (259 of 754) cite a value that sits on the country's official
ladder --- a superseded value or a future stage --- rather than a fabricated
one (Section~\ref{sec:results:t1}). The floor is real, shallower than first
reported, and concentrated in normative recall (T1).

\begin{table}[h]
\centering
\caption{Composite accuracy by task type (25 countries; T1 on 24, Egypt having
no key). T1, T2 and T3 are code-adjudicated where the returned value is
checkable against the official register; T4 is panel-scored; T5 retains the
collection judge's scores.}
\label{tab:conf-task}
\footnotesize
\setlength{\tabcolsep}{8pt}
\renewcommand{\arraystretch}{1.15}
\begin{tabular}{lrr}
\toprule
Task & $N$ & Mean \\
\midrule
T5 (applied recommendation)  & 1{,}320 & 0.769 \\
T3 (health-evidence synth.)  & 1{,}329 & 0.548 \\
T4 (policy instruments)      & 1{,}310 & 0.520 \\
T2 (local factual datum)     & 1{,}335 & 0.473 \\
T1 (technical standard)      & 1{,}279 & 0.394 \\
\bottomrule
\end{tabular}
\end{table}

\subsection{H2 --- asking in the local language lowers accuracy}
\label{sec:results:h2}

For the nine countries with a target native language, each prompt is rendered in
both English and the native language with content held fixed, so the contrast is
within-country and within-model. Replicates are averaged within each (model,
prompt) cell before pairing, since pairing replicate-to-replicate would inflate
$n$ and make the test anticonservative. Across the 839 matched cells,
native-language prompting \emph{lowers} mean composite accuracy by 5.0 percentage
points (Wilcoxon signed-rank $p=5\times10^{-16}$). The penalty tracks the resource
level of the language, and all three are significant (Table~\ref{tab:h2}): Hindi
$-11.4$~pp, Spanish $-4.5$~pp, Portuguese $-4.2$~pp. India, which scores highest
of all 25 countries in English, drops the furthest in Hindi.

Because H2 carries the study's central claim, we tried to make it disappear rather
than to confirm it, and it held under every attack: it is not driven by any one
country (leave-one-out $-4.4$ to $-5.5$~pp, all significant), by any one model
($-4.5$ to $-5.2$~pp), or by outlying cells (5\% trimmed $-4.5$~pp,
$p=6\times10^{-20}$); it survives every weighting of the composite; and it appears
in every task and under both persona conditions.

Two attacks answer the objections a sceptical reader raises first. The first is
that language-model judges might penalise non-English prose for style rather than
content, so we test a binary outcome no judge touches: does the response contain a
value the extractor can compare with the register? The extractor reads Devanagari
numerals, decimal commas and units written out in words, so a Hindi or a
Spanish-language number is not penalised for its notation
(Section~\ref{sec:methods}). English prompts yield a verifiable value $71.9\%$ of
the time and native-language prompts $62.5\%$, a paired difference of $-9.4$~pp,
with the native version worse in 134 of the 218 discordant pairs (sign test
$p=9\times10^{-4}$): Hindi $-39.5$~pp, Spanish $-13.4$~pp, while Portuguese, the
best-resourced of the three, shows no such loss ($+3.5$~pp). Asking in the local
language does not merely lower a judged score; where the language is scarce it
makes the model substantially less likely to return a checkable number at all.

The second is that the native prompts, produced by an LLM translator, might simply
ask something different. We tested this on all 90 unique native prompts --- a
census, not a sample --- back-translating each with a \emph{different} model and
having two further models audit the result against the English original. On the
two items that could confound the effect, fidelity is complete: no prompt changed
the country and none changed the requested quantity ($90/90$ on both). Restricting
H2 to the prompts verified as faithful leaves it essentially unchanged
($-4.8$~pp, $p=2\times10^{-8}$) against $-5.2$~pp among the divergent ones, a
difference not distinguishable from zero (permutation $p=0.80$). Translation
fidelity does not predict the size of the penalty, which is what one expects if
translation is not causing it. Separating the pairs by the instrument that scored
the native cell, the penalty is $-3.8$~pp where code decided ($n=333$,
$p=0.004$), $-8.6$~pp where the panel decided ($n=314$, $p<10^{-21}$) and
$-1.1$~pp, not significant, where the collection judge's score remains ($n=192$,
essentially T5, the open recommendation with no register value). The
code-adjudicated subset is the smaller of the two because code acts only where a
number exists, so it conditions on the model having produced a value in
\emph{both} languages and selects away the severest failure mode; the panel
subset carries that failure mode, which is why judge leniency cannot be assumed
to cancel across languages and why the judge-free outcome above is the one that
settles the point.

The applied implication is the opposite of the intuition. A manager who suspects
an English-language model serves their country poorly, and responds by asking in
the national language, makes the answer \emph{worse} --- and worse by the largest
margin precisely where the language is least represented in training corpora.

\begin{table}[h]
\centering
\caption{H2: native-minus-English accuracy by native language (paired within
country and model, replicates averaged within cell, 25-country sample).}
\label{tab:h2}
\footnotesize
\setlength{\tabcolsep}{8pt}
\renewcommand{\arraystretch}{1.15}
\begin{tabular}{llrrr}
\toprule
Native language (Joshi class) & Countries & $\Delta$ (pp) & $p$ & $N$ cells \\
\midrule
Spanish (5)    & MEX, ARG, PER, COL, CHL & $-4.5$  & $2\times10^{-10}$ & 457 \\
Portuguese (4) & BRA, PRT, AGO           & $-4.2$  & $8\times10^{-5}$  & 311 \\
Hindi (4)      & IND                     & $-11.4$ & $4\times10^{-4}$  & 71 \\
\midrule
\multicolumn{2}{l}{Overall}              & $-5.0$  & $5\times10^{-16}$ & 839 \\
\bottomrule
\end{tabular}
\end{table}

\subsection{H6 --- persona does not narrow the geographic gap}
\label{sec:results:h6}

The persona factor was fully crossed, so H6 is tested as a
difference-in-differences. The public-environmental-manager frame leaves the
North/South gap essentially unchanged: $+4.8$~pp under the neutral frame versus
$+5.2$~pp under the manager frame, a DiD of $+0.5$~pp, far below the
pre-specified 5~pp threshold; a country-level permutation test (5{,}000
permutations of the tier label) gives $p=0.29$. H6 is not supported:
presenting the query as coming from a public environmental manager neither
narrows nor amplifies the Global South disadvantage, consistent with evidence
that expert personas do not reliably improve factual
accuracy~\citep{zheng2024helpful,mollick2025personas}. A manipulation check finds
that $50.2\%$ of persona-condition responses explicitly acknowledge the role
frame, so the null is not an artefact of the frame being universally ignored.

\subsection{H1 --- a tier gap that holds, a gradient that does not}
\label{sec:results:h1}

H1 was pre-specified in two parts, and they separate cleanly.

The ten Global North countries average $0.572$ against
$0.521$ for the fifteen Global South countries, a gap of $+5.0$ percentage points
whose 95\% bootstrap confidence interval over countries is $[+1.5,+8.4]$~pp and
\emph{excludes zero} (permutation of the tier label, country as unit,
$p=0.021$). It is stable to dropping any single country (leave-one-out
$[+4.4,+5.9]$~pp) and to the composite weighting ($+4.2$~pp under equal weights,
$+8.7$~pp on the factual subcomponent, where it widens). Because seven of the ten
Global North countries entered with the extension, we also report the gap where
the wave cannot explain it: within the extension alone (seven Global North
against Colombia, Chile and Angola) it is $+5.9$~pp $[+2.3,+9.6]$, and on the
pre-specified fifteen (three Global North against twelve) $+6.0$~pp
$[+1.7,+10.2]$, a point estimate of the same size on a sample too small to make
it significant by permutation ($p=0.13$). The gap is a property of the data, not
of the wave that collected it.

\emph{Where the gap lives, and why the composite understates it.} The composite
averages five tasks and the gap is not spread evenly across them. Re-estimating
the same contrast as a binomial mixed model on each task separately --- outcome
coded as returning the register value, with crossed random intercepts for country
and model --- shows the disadvantage scaling with how much the task depends on a
fact specific to that country (Table~\ref{tab:tier-by-task}). On the binding
national standard (T1), scored by code against the instrument in force, the
Global South odds of returning the register value are $0.33$ of the Global
North's ($50.7\%$ versus $31.6\%$ raw), and on policy instruments (T4) also
$0.33$; on health-evidence synthesis (T3), where the answer is an international
estimate rather than a national one, the ratio rises to $0.69$; and on open
recommendation (T5), the one task with no register value to get right, the gap
disappears entirely. The T1 ratio survives the two objections its composition
invites (Section~\ref{sec:results:t1}): setting aside the four Global South
countries whose key is not a plainly verified national value it is $0.39$
$[0.31,0.50]$, and on the pre-specified fifteen it is $0.27$ $[0.21,0.34]$. The
gap is not weak --- it is \emph{concentrated}, and the composite dilutes it by
averaging tasks where it is severe with a task where it does not exist. On the
question a public manager most often brings to these systems, what is the
standard in force here, a Global South country is served at a third of the odds
of a Global North one.

\begin{table}[h]
\centering
\caption{Global South versus Global North odds of a factually correct answer,
by task. The outcome is factual accuracy $\ge 0.5$: the code verdict where the
task is adjudicated by code (all of T1; the resolvable part of T2 and T3), the
panel or judge score otherwise, so for T5, which has no register value, it is the
judge's factual rating. Binomial mixed model, crossed random intercepts for
country and model.}
\label{tab:tier-by-task}
\footnotesize
\setlength{\tabcolsep}{8pt}
\renewcommand{\arraystretch}{1.15}
\begin{tabular}{llrrrl}
\toprule
Task & Answer is & GN & GS & OR & 95\% interval \\
\midrule
T1 (technical standard)     & a national value        & 50.7\% & 31.6\% & 0.33 & $[0.27,0.40]$ \\
T4 (policy instruments)     & national instruments    & 62.0\% & 44.7\% & 0.33 & $[0.26,0.40]$ \\
T2 (local measured datum)   & a national value        & 50.7\% & 37.8\% & 0.47 & $[0.39,0.55]$ \\
T3 (health-evidence synth.) & an international estimate& 59.1\% & 52.3\% & 0.69 & $[0.59,0.80]$ \\
T5 (applied recommendation) & no register value       & 99.6\% & 99.7\% & 3.58 & $[0.84,15.17]$ \\
\bottomrule
\end{tabular}
\end{table}

At $n=25$, Spearman
$\rho(\text{accuracy},\text{HDI})=+0.36$ does not reach conventional significance
($p=0.076$; Mann-Kendall $p=0.065$) and stays well below the pre-specified
$\rho\geq0.55$ threshold, which is the criterion H1 was fixed against. At the
pre-specified $n=15$ it is essentially absent ($\rho=+0.08$), and no
leave-one-country-out re-estimate reaches $0.55$ either (range $[+0.30,+0.48]$);
under equal weights it is $0.35$ and under PCA weights $0.37$. We
therefore report a tier difference we can demonstrate and a development gradient
we cannot: the disadvantage is established as a difference between groups, while
its shape across the development range remains exploratory. Under Bonferroni-Holm
across the primary family $\{$H1-gradient, H4, H6$\}$, no test rejects its null.

The country ordering retains large axis-orthogonal heterogeneity (Supplementary
Material): India ($0.651$, Global South) tops all countries, Indonesia and
South Africa outrank most of the Global North, and Canada and France sit low. The
disadvantage is a tier difference layered over substantial country-specific
variation, and that heterogeneity is itself part of why the gradient does not
resolve.

\subsection{H4 --- corpus representation: country coverage, inseparable from development}
\label{sec:results:h4}

The pre-specified mechanism was training-corpus representation, proxied by
Wikipedia article and edition count. That proxy is null: Spearman $\rho=+0.05$
($p=0.82$) uncontrolled, and still null partialling out HDI and log GDP per
capita. A post-hoc country-coverage proxy reaches $\rho=+0.37$ ($p=0.070$) and
attenuates after adjusting for HDI (partial $\rho=+0.19$, $p=0.37$). At country
level neither channel is established, because with 25 countries coverage and
development are collinear.

\subsubsection*{Changing the unit of variation resolves what the country level
cannot}

The obstacle is confounding rather than absence of signal, and the way out does
not require more countries. Each country answers five tasks that differ in how much
the answer depends on a fact about \emph{that} country: T1, T2 and T4 do; T3 (an
international WHO estimate) and T5 (a generic recommendation) do not. Countries
score $0.462$ on the first block and $0.658$ on the second, a specificity deficit
of $0.195$. If what limits the model is how much the corpus says about a given
country, that deficit should be smaller where coverage is greater. Estimating the
interaction between coverage and task dependence at the response level, with a
fixed effect for country, gives $\beta=+0.030$ (SE $0.005$,
$p=9\times10^{-10}$, $n=6{,}573$). The fixed effect absorbs every country-level
\emph{main} effect --- development, income, official language --- but not their
interaction with task dependence, and that is where the rival lives: a richer
country may simply have more of its regulation online, feeding both the model's
training data and the ease of building the key.

Four checks separate the coverage signal from artefact (Table~\ref{tab:h4-within}).
The \emph{purest} contrast --- T1, wholly national, against T5, wholly generic ---
yields the \emph{largest} coefficient; dropping T5 or T4 leaves it intact; a
placebo inverting the labels returns the exact mirror image; and entering the
language-corpus proxy alongside it, coverage survives ($\beta=+0.030$,
$p=2\times10^{-9}$) while language does not ($\beta=+0.004$, $p=0.45$). Because
every response here was elicited in English, that last result was predicted and
is the falsifiable part of the design. The rival it does not eliminate is
development. Coverage and HDI correlate at $\rho=0.65$ across the 25 countries,
and when the two interactions enter the same model the ordering reverses:
HDI$\times$dependence survives ($\beta=+0.027$, SE $0.007$, $p=5\times10^{-5}$)
and coverage$\times$dependence does not ($\beta=+0.012$, SE $0.007$,
$p=0.063$). Within the nine anglophone countries, where the language of the
question is the same by construction, coverage predicts the deficit at the full
magnitude ($\beta=+0.024$, $p=0.003$); that removes language as the explanation,
not development.

The resulting claim is therefore narrower than ``coverage causes the gap'', and
narrower than we first wrote it. What the within-country design establishes is
that the country-specific deficit is smaller where the world records more about
the country --- whether that record is measured as encyclopedic coverage or as
development, which in this sample are two readings of the same variable. It rules
out the language-corpus channel; it does not pick coverage over development, and
we report the mechanism as identified only up to that pair.

\begin{table}[h]
\centering
\caption{Country coverage $\times$ task dependence, response-level mixed model
with a fixed effect for country. Positive $\beta$ = greater coverage, smaller
deficit on country-dependent tasks.}
\label{tab:h4-within}
\footnotesize
\setlength{\tabcolsep}{8pt}
\renewcommand{\arraystretch}{1.15}
\begin{tabular}{lrl}
\toprule
Task classification & $\beta$ & $p$ \\
\midrule
T1, T2, T4 versus T3, T5 (primary) & $+0.030$ & $9\times10^{-10}$ \\
T1 versus T5 (purest contrast)     & $+0.045$ & $2\times10^{-11}$ \\
Excluding T5 (ceiling at 99.7\%)   & $+0.026$ & $3\times10^{-5}$ \\
Excluding T4 (most interpretive)   & $+0.034$ & $1\times10^{-9}$ \\
Placebo: labels inverted           & $-0.030$ & $9\times10^{-10}$ \\
With HDI$\times$dependence: coverage & $+0.012$ & $0.063$ \\
\quad HDI                          & $+0.027$ & $5\times10^{-5}$ \\
\bottomrule
\end{tabular}
\end{table}

The corpus account shows itself more cleanly in H2, where the unit is the language: the
per-language penalty falls monotonically with mC4 token
volume~\citep{xue2021mt5} and OSCAR size~\citep{abadji2022oscar}, with Hindi, the
smallest corpus, carrying the largest penalty --- though with three languages this
remains descriptive.

\subsection{H5 --- open versus closed-accessible}
\label{sec:results:h5}

Closed-accessible models average $+13.0$~pp above open-weight ones, against the
optimistic H5 framing. The open arm includes larger models, so the gap is not
attributable to scale alone, although DeepSeek-V3 ranks first overall. H5 is
interpreted descriptively.

\subsection{Ground-truth provenance, and what the T1 key decides}
\label{sec:results:gt}
\label{sec:results:t1}

Every key comes from an official register, anchored by documentary verification
and reproducible from the cited URLs: T1 from the national regulation, T2 from WHO
Ambient Air Quality Database v6.1, T3 from WHO GHO indicator AIR\_41, T4 from the
UNEP Global Assessment of Air Pollution Legislation. Three countries have \emph{no}
national PM\textsubscript{2.5} standard, confirmed from the official text ---
Nigeria regulates PM\textsubscript{10} only, Argentina has no binding federal
value, Angola no national legislation --- which is where a faithful model must
decline to state a non-existent threshold. Two keys are not fully verified
(Bangladesh, scored against its 2005 gazette; Egypt, with no retrievable value
and therefore no T1 key), and all five of these cases are Global South, so the
composition of the key is itself a candidate explanation for the T1 gap. It is
not: set them aside and the Global South odds on T1 are $0.39$ $[0.31,0.50]$ of
the Global North's, against $0.33$ with them. The full registry is in
Supplementary Material. Wikipedia and Wikidata appear only as the explanatory
variable of the mechanism test; no answer is scored against them.

\emph{The key in force, and the ladder around it.} Several instruments set a
ladder of values rather than one: Brazil's CONAMA 506/2024 runs 20, 17, 15, 10,
5~$\mu$g/m\textsuperscript{3} across five stages, of which 17 has been in force
since January 2025; the European limit is 25 with 10 from 2030; the United
Kingdom's limit value is 20 with an English target of 10 by 2040; the United
States moved from 12.0 to 9.0 in 2024. The T1 key is the value in force in the
collection window, and it is exacting: no model returns 17 for Brazil (the modal
answer is 15, the next stage), almost none returns 20 for the United Kingdom
(most return the 2040 target or the superseded 25), and fewer than a quarter
return 9.0 for the United States. Of the 754 wrong T1 answers, 259 cite a value
that sits somewhere on the country's official ladder --- a stage not yet in
force or a value since replaced --- which is confusion about \emph{which} value
binds, not fabrication, and is the failure a training cut-off would produce. The
distinction does not rescue the tier comparison, it sharpens it: scoring any
ladder value as correct raises Global North accuracy on T1 from $0.51$ to $0.77$
and Global South accuracy only from $0.32$ to $0.40$, because the staged
instruments are mostly in the North, and the odds ratio falls to $0.12$
$[0.10,0.15]$. Where the register offers several plausible numbers the models
know them; where it offers one, they often do not.

\subsection{Qualitative failure modes}
\label{sec:results:qual}

Three failure modes recur behind the composite, with verbatim examples in
Supplementary Material. Models \emph{fabricate standards that do not exist} for
the three countries with no national norm: in 168 of the 256 scored T1 cells for
Angola, Argentina and Nigeria the answer asserts a specific national value, in
33 it correctly states that none exists, and in 31 it declines to answer, and the
invented values disagree among replicates; they \emph{miss the binding value},
answering 15~$\mu$g/m\textsuperscript{3} for Brazil where the stage in force is
17 (CONAMA 506/2024), a confusion of stage rather than an invention; and they
\emph{degrade in the native language}, as when GPT-OSS~120B returned India's
correct 40~$\mu$g/m\textsuperscript{3} in English and an \emph{empty} answer in
Hindi.

\subsection{Inter-judge reliability}
\label{sec:results:reliability}

Reliability is reported on the sample that produces the effects, not on a
calibration subset. All 3{,}190 items requiring a judge were scored by all three
panel members, giving Krippendorff's $\alpha=0.527$, single-judge
ICC$(2,1)=0.558$, and --- the operative quantity, since the analysis uses the panel
mean --- ICC$(2,3)=0.791$.

The panel also isolates a design rule that generalises beyond this study. Where
the judging
prompt asked the model to compare a returned concentration against a tolerance
band expressed in words --- that is, to perform the arithmetic itself ---
inter-judge concordance on those items was $r=-0.04$. Pre-computing the admissible
range and asking only whether the value falls inside it raised concordance on the
same items to $r=+1.00$, and the 274 scores collected under the ambiguous wording
were discarded and recollected. Apparent judge disagreement can be an artefact of
what the judge is asked to compute. Arithmetic belongs to code.

\subsection{Robustness}
\label{sec:results:robust}

The two supported effects were tested against seven attacks and survive all of
them; the development gradient stays below the pre-specified $0.55$ under every
one of them (its highest value, $0.48$, is a leave-one-out re-estimate), which is
what makes its exploratory status a finding rather than an omission. Full outputs
are in Supplementary Material.

\emph{No single country, and no weighting, carries the result.}
Leave-one-country-out re-estimation holds the tier gap in $[+4.4,+5.9]$~pp and the
native-language penalty at $-4.4$ to $-5.5$~pp. Under equal weights the gap is
$+4.2$~pp and the language penalty $-4.8$~pp; restricting to factual accuracy
against the register value widens them to $+8.7$ and $-6.1$~pp and deepens the
recall floor, so both effects grow where the measurement is most objective. The
two replicates of each cell differ by $8.5$~pp on average (within-cell SD
$10.1$~pp), so the $-5.0$~pp language penalty, averaged over 839 pairs, stands
roughly ten standard errors from the noise two replicates leave unresolved.

\emph{The tier gap survives a distribution-free test and two model families.} A
permutation test shuffles the North/South labels among the 25 countries ten
thousand times and counts how often chance alone produces a gap this large; it
assumes nothing about the distribution and treats the country, not the answer, as
the independent unit. Chance clears the observed gap $2.1\%$ of the time
($p=0.021$). A Bayesian hierarchical model with crossed random intercepts agrees
($\beta=-0.048$, 94\% HDI $[-0.090,-0.009]$, $P(\beta<0)=0.987$), while a
frequentist mixed model returns the same magnitude and sign but leaves it marginal
($\beta=-0.063$, $p=0.087$), because it charges the large between-country
heterogeneity to the error term. Two of three tests clear zero and all three agree
on magnitude, so we report the gap as supported and specification-sensitive, with
an E-value of $1.70$ at the point estimate and $1.30$ at the confidence limit
nearest the null~\citep{vanderweele2017evalue} --- a country-level confounder of
modest strength would suffice to explain it away.

\emph{The gap tracks world-system position, not national wealth.} Re-estimated
under three development-based partitions of the same 25 countries, the gap is
significant under UNCTAD and non-significant under every split drawn on HDI
(Supplementary Material). This coheres with the rest of the paper: the HDI
gradient is itself null, so a split on HDI should not be expected to produce a
gap, and that UNCTAD separates the data where HDI does not points to position in
the world system rather than income as the operative
variable~\citep{quijano2000coloniality,mohamed2020decolonial}. The reading is post
hoc, so we claim the gap for the North/South distinction as conventionally drawn,
not as a general statement about development.

\emph{What the null gradient does and does not exclude.} Simulating
25-country samples at known correlations, the design has $77\%$ power at
$\rho=0.55$ --- the threshold fixed in advance --- and a minimum detectable effect
of $\rho=0.56$. A gradient of the magnitude we committed to calling substantial is
excluded; a weak one is not (power $48\%$ at $\rho=0.40$). The ten countries added
post hoc fall \emph{within} the HDI range the original fifteen already spanned, so
the extension raises the density of observations rather than the range of the
predictor.

\subsection{Summary}
\label{sec:results:summary}

Table~\ref{tab:hyp-summary} maps each hypothesis to its finding, the candidate
mechanism, and an explicit evidence status, so no claim is read as stronger than
the data support. Three effects are supported and survive every robustness attack
in Section~\ref{sec:results:robust}; two pre-specified expectations are not met;
and the corpus mechanism is identified within countries only up to the pair
coverage-or-development, which this design cannot separate.

\begin{table}[!ht]
\centering
\caption{Hypotheses, findings, candidate mechanisms, and evidence status. Grouped
by what the evidence supports, not by hypothesis number.}
\label{tab:hyp-summary}
\footnotesize
\setlength{\tabcolsep}{8pt}
\renewcommand{\arraystretch}{1.15}
\renewcommand{\arraystretch}{1.15}
\begin{tabular}{@{}>{\raggedright\arraybackslash}p{0.235\linewidth}%
                  >{\raggedright\arraybackslash}p{0.315\linewidth}%
                  >{\raggedright\arraybackslash}p{0.235\linewidth}@{}}
\toprule
Hypothesis & Finding & Mechanism \\
\midrule
\multicolumn{3}{@{}l}{\emph{Supported}} \\
\addlinespace[2pt]
H2 language        & $-5.0$~pp; Hindi $-11.4$    & language corpus \\
H3 regional model  & $\delta=-0.48$, worst of 14 & scale over tuning \\
H1 tier gap        & $+5.0$~pp; OR $0.33$ on T1  & tier \\
\midrule
\multicolumn{3}{@{}l}{\emph{Supported, mechanism not separable}} \\
\addlinespace[2pt]
H4 within country  & $\beta=+0.030$ alone; $+0.012$ ($p=0.063$) with HDI & coverage or development \\
\midrule
\multicolumn{3}{@{}l}{\emph{Not supported}} \\
\addlinespace[2pt]
H1 gradient        & $\rho=0.36$, below the $0.55$ threshold & development \\
H4 pre-specified   & language-corpus proxy null  & language corpus \\
H6 persona         & DiD $+0.5$~pp, $p=0.29$     & --- \\
\midrule
\multicolumn{3}{@{}l}{\emph{Descriptive only}} \\
\addlinespace[2pt]
H5 open/closed     & $+13.0$~pp for closed       & access type \\
\bottomrule
\end{tabular}

\vspace{4pt}
\begin{minipage}{0.75\linewidth}
\footnotesize\emph{Note.} H2's mechanism rests on three native languages and is
descriptive. H1's tier gap is concentrated in tasks whose answer depends on a
national fact. H4 within country rules out the language-corpus channel but
cannot separate coverage from development, and is not causal identification.
\end{minipage}
\end{table}